% Generated draft: requires mathematical and source review.
\documentclass[11pt]{article}
\usepackage[margin=25mm]{geometry}
\usepackage{fontspec}
\setmainfont{Latin Modern Roman}
\usepackage{amsmath,amssymb,mathtools,unicode-math}
\setmathfont{Latin Modern Math}
\usepackage{enumitem,needspace,xurl,hyperref}
\setlength{\parindent}{0pt}
\setlength{\parskip}{5pt}
\setlength{\emergencystretch}{3em}
\newcommand{\N}{\mathbb{N}}
\newcommand{\Z}{\mathbb{Z}}
\newcommand{\Q}{\mathbb{Q}}
\newcommand{\R}{\mathbb{R}}
\newcommand{\C}{\mathbb{C}}
\newcommand{\F}{\mathbb{F}}
\newcommand{\A}{\mathbb{A}}
\newcommand{\PP}{\mathbb{P}}
\newcommand{\E}{\mathbb{E}}
\newcommand{\eps}{\varepsilon}
\newcommand{\dd}{\,\mathrm{d}}
\begin{document}
\section*{UnsolvedMath ID: 8800057}
\textbf{Source problem code:} catalogue entry
\subsection{Definitions and mathematical statement}
Pairing-based assumptions concern parameterized groups $G_1,G_2,G_T$ with an efficiently computable bilinear map $e:G_1\times G_2\to G_T$, together with a specified challenge distribution and required output or decision predicate. For a security parameter $\kappa$, an asymptotic hardness claim must bound the success probability of algorithms running in the allowed resource class; a negligible function is eventually smaller than $\kappa^{-c}$ for every fixed $c>0$.

The source asks to break or determine the actual hardness of the ``weaker'' assumptions used in pairing cryptography. It does not specify the challenge tuples, group family, pairing type, distribution, computational versus decisional task, parameter growth, adversary model, or baseline assumption relative to which weaker is meant. These data are needed before success probabilities, reductions or lower bounds can be stated.

The summary mentions $k$-bilinear Diffie--Hellman inversion and co-BDH, but supplies no definitions or choice among their variants. Their names do not determine a unique mathematical problem here. Concrete algorithms, reductions between assumptions, and generic-group-model lower bounds are separate kinds of hardness evidence. A generic-model result restricts the available representation/oracles and is not automatically a standard-model result for every instantiated group. No specific assumption, reduction direction or claim of unconditional security is added to this programme.
\subsection{Short English statement}
Assess the source's weaker pairing-based security assumptions through precise algorithms or hardness comparisons. The exact assumptions, distributions, group families and computational model must first be specified.
\subsection{Sources}
\begin{itemize}
\item Local UnsolvedMath registry, permanent ID 8800057, exact source code catalogue entry, ``Breaking weaker pairing assumptions''; attributed to Tanja Lange, no proposed year supplied. Registry status open, classification OPEN-TRIAGE.
\item \emph{IPAM Workshop on Number Theory and Cryptography---Open Problems} (2006), attachment 6239, PDF page 41, printed slide 39, \url{https://www.ipam.ucla.edu/programs/workshops/workshop-i-number-theory-and-cryptography-open-problems/?tab=schedule}. Supplied pdf-and-legacy-ppt visual audit accessed July 29, 2026; provenance records 17 PDFs/two PPT decks inspected with repeated builds excluded. No PDF/PPT was accessed here.
\item The supplied summary names $k$-bilinear Diffie--Hellman inversion, co-BDH and related assumptions, says their relationships to standard assumptions are unresolved in its assessment, and reports no known break. It supplies no exact challenges, parameter ranges, theorem references or dated attack coverage; those broad status claims are not independently verified.
\item Provenance flags current primary-literature status and redistribution terms NEEDS\_REVIEW.
\end{itemize}
\end{document}
